\subsection{Constructing and resolving an unstable context}
Let $U=\operatorname{Top}_k(P(Q)\cup F)\setminus(H(Q)\cup F)$.
These claims have unresolved support and enter the possible context.
Their status determines whether date uncertainty can change retrieval.

\begin{theorem}[Constructive instability and refinement]
\label{thm:refinement}
If $U\ne\varnothing$, we construct two complete date assignments with different contexts in $O(n)$ time.
Here, $n$ counts the event dates.
Every date refinement that stabilizes the context makes each claim in $U$ guaranteed or impossible.
\end{theorem}

Each $c\in U$ has fewer than $k$ possible claims above it.
Any timeline supporting $c$ therefore selects it.
We place every accepted replacement no later than $c$ or after a supporting query time.
A second timeline places $c$ after the query, or a direct witness by the query start.
Date refinement can only remove possible claims.
Thus a remaining possible pivot cannot leave the possible context through rank changes.
It must become guaranteed.
The assignments use only event bounds and the compact certificate.
Let $N$ count ranked excerpts, including fallback text.
Full context replay takes $O(n+m+N)$ time with the accepted pair list.
Appendix~\ref{app:refinement} gives the construction and proof.

The set $U$ identifies support decisions that every stabilizing refinement must settle.
New unresolved claims can enter after current pivots disappear.
We recompute the frontier after each refinement.

\subsection{A tractability boundary}
Possible support concerns one claim at a time.
Possible context membership also depends on which higher-ranked claims coexist.
We ask whether a given claim enters the context under \emph{any} allowed timeline.
\begin{theorem}[Possible context membership]
\label{thm:hardness}
When $k$ is part of the input, deciding possible context membership is NP-hard.
\end{theorem}
This holds even for a point query and a bipartite gate with one edge layer.
\begin{proof}
Take a Set Cover instance with elements $E$, sets $S_1,\ldots,S_s$, and budget $1\le B\le s$~\citep{karp1972reducibility}.
Each set becomes a claim $v_j$ with start bounds $[1,3]$.
Create $B+1$ copies of each element, all starting at time 0.
Add a candidate $c$ starting at time 0.
A pair $v_j\to e$ targets every copy of each element in $S_j$.
No other pairs are accepted.
Rank every set claim and element copy above $c$.
Use point query 2 and context budget $B+1$.

For a cover of at most $B$ sets, start its set claims at time 1.
Start the others at time 3.
Every element copy retires before the query.
At most $B$ higher-ranked set claims remain active, so $c$ enters the context.
Conversely, a context containing $c$ has at most $B$ active higher-ranked claims.
Its active set claims must cover every element.
An uncovered element would leave $B+1$ active copies above $c$ and exclude it.
Thus these claims define a cover of size at most $B$.
The construction is polynomial because $B\le s$.
\end{proof}
Appendix~\ref{app:refinement} also gives a reduction where every start precedes or equals the query time.
Our stability test instead needs one unresolved claim inside the possible prefix.
Its rank guarantees selection whenever it has support.
The test needs only individual support, even when shared events couple claim eligibility.
This makes exact stability and its witnessing timelines tractable without computing every attainable context member.
